% 完整翻译：Coordinate Systems（坐标系）
\section{坐标系}

本节完整翻译并重述自参考资料中关于“坐标系”的段落，目的是以中文系统地介绍用于向量分析的常见坐标表示、微元与几何关系，以及不同坐标系之间的相互转换与应用。我们着重介绍直角（笛卡尔）、圆柱和球坐标系。

\subsection{坐标系的定义与选取}
坐标系是一种以参考原点为基准，唯一标识空间中任意位置的方法。任一点的位置可由三条两两互相垂直的曲面（在本教材中通常取三平面）的交线来确定；在该点处，这三条曲面的法向量定义了坐标轴的方向。虽然物理问题的结果不依赖于所选坐标系，但通过利用问题的对称性，恰当选择坐标系可以显著简化计算。在本书中，我们仅采用直角（Cartesian）、圆柱（cylindrical）和球（spherical）三类常用坐标系。

\subsection{直角坐标系（Rectangular/Cartesian）}
直角坐标由三条两两垂直的平面的交线定义，通常记作 $(x,y,z)$。对应的坐标轴在任一点处分别垂直于恒定 $x$, $y$, $z$ 的平面（见图 1-1）。选定原点为 $(0,0,0)$ 后，空间中任一点可用沿 $x,y,z$ 方向从原点到该点的距离来表示。

按约定我们使用右手坐标系：把右手四指由 $x$ 指向 $y$ 弯曲，拇指所指方向为 $z$。这一约定用于消除方向上的模糊。直角坐标的单位向量通常记为 $\mathbf{i}_x,\mathbf{i}_y,\mathbf{i}_z$（或简写 $\mathbf{i},\mathbf{j},\mathbf{k}$），它们在空间中处处方向相同，因此直角坐标在许多情形下是最简便的选择。

微元：在点 $(x,y,z)$ 处沿三个坐标方向的微增量为 $\mathrm{d}x,\mathrm{d}y,\mathrm{d}z$，从而微小体积为
\[
\mathrm{d}V=\mathrm{d}x\,\mathrm{d}y\,\mathrm{d}z.
\]
相应的面元按垂直的坐标方向下标记（例如垂直于 $x$ 的面元记作 $\mathrm{d}S_x$）。

\subsection{圆柱坐标系（Circular cylindrical）}
当问题关于一条直线（取为 $z$ 轴）具有轴对称性时，圆柱坐标系 $(r,\phi,z)$ 很方便。一个点由以下三个表面交给出：半径为 $r$ 的圆柱面、恒定 $z$ 的平面，以及与 $x$ 轴夹角为 $\phi$ 的平面。直角坐标与圆柱坐标的正交投影关系为
\[
x=r\cos\phi,\qquad y=r\sin\phi,\qquad z=z.
\]

对应的自然基向量（未单位化）为 $\mathbf{e}_r,\mathbf{e}_\phi,\mathbf{e}_z$，其中 $\mathbf{e}_r$ 与 $\mathbf{e}_\phi$ 的方向随角 $\phi$ 而改变，而 $\mathbf{e}_z$ 方向不变。尺度因子（scale factors）为 $h_r=1,\ h_\phi=r,\ h_z=1$，因此线元为
\[
\mathrm{d}s^2=\mathrm{d}r^2+r^2\mathrm{d}\phi^2+\mathrm{d}z^2,
\]
面元和体元中会出现因子 $r$，例如体元为
\[
\mathrm{d}V=r\,\mathrm{d}r\,\mathrm{d}\phi\,\mathrm{d}z.
\]

圆柱坐标下的单位向量 $\hat{\mathbf{e}}_r$ 与 $\hat{\mathbf{e}}_\phi$ 可通过直角坐标的单位向量表示，且满足右手坐标系的方向约定（$\hat{\mathbf{e}}_r\times\hat{\mathbf{e}}_\phi=\hat{\mathbf{e}}_z$）。

\subsection{球坐标系（Spherical）}
当问题围绕某点有球对称性时，球坐标系 $(r,\theta,\phi)$ 更为合适。这里 $r$ 为径向距离，$\theta$ 为极角（从 $z$ 轴向外量起），$\phi$ 为方位角（绕 $z$ 轴）。直角坐标与球坐标的关系为
\[
x=r\sin\theta\cos\phi,\qquad y=r\sin\theta\sin\phi,\qquad z=r\cos\theta.
\]

在球坐标中尺度因子为 $h_r=1,\ h_\theta=r,\ h_\phi=r\sin\theta$，因此线元为
\[
\mathrm{d}s^2=\mathrm{d}r^2+r^2\mathrm{d}\theta^2+r^2\sin^2\theta\,\mathrm{d}\phi^2,
\]
体元为
\[
\mathrm{d}V=r^2\sin\theta\,\mathrm{d}r\,\mathrm{d}\theta\,\mathrm{d}\phi.
\]

对应的单位向量 $\hat{\mathbf{e}}_r,\hat{\mathbf{e}}_\theta,\hat{\mathbf{e}}_\phi$ 在空间中随位置改变；三者构成右手系统。

\subsection{各坐标系的微元汇总}
常用的微元长度、面元与体元在三类坐标系中的形式可归纳如下（与教材表格对应，省略表格排版但给出常用结果）：
- 笛卡尔：$\mathrm{d}l_x=\mathrm{d}x,\ \mathrm{d}S_x=\mathrm{d}y\,\mathrm{d}z,\ \mathrm{d}V=\mathrm{d}x\,\mathrm{d}y\,\mathrm{d}z$。
- 圆柱：$\mathrm{d}l_r=\mathrm{d}r,\ \mathrm{d}l_\phi=r\,\mathrm{d}\phi,\ \mathrm{d}S_r=r\,\mathrm{d}\phi\,\mathrm{d}z,\ \mathrm{d}V=r\,\mathrm{d}r\,\mathrm{d}\phi\,\mathrm{d}z$。
- 球坐标：$\mathrm{d}l_r=\mathrm{d}r,\ \mathrm{d}l_\theta=r\,\mathrm{d}\theta,\ \mathrm{d}l_\phi=r\sin\theta\,\mathrm{d}\phi,\ \mathrm{d}V=r^2\sin\theta\,\mathrm{d}r\,\mathrm{d}\theta\,\mathrm{d}\phi$。

\subsection{坐标和分量之间的转换}
若需把向量或位置从一种坐标表示转换到另一种，可使用几何关系表（如教材中的表 1-2）。例如圆柱与直角之间的转换为
\[
r=\sqrt{x^2+y^2},\qquad \phi=\arctan\frac{y}{x},\qquad z=z,
\]
球坐标与直角之间的转换为
\[
r=\sqrt{x^2+y^2+z^2},\qquad \theta=\arccos\frac{z}{r},\qquad \phi=\arctan\frac{y}{x}.
\]

基向量的变换遵循链式法则；若从坐标 $x^i$ 变换到新坐标 $\bar x^a=\bar x^a(x^i)$，则坐标基底的变换满足
\[
\frac{\partial}{\partial \bar x^a}=\frac{\partial x^i}{\partial \bar x^a}\frac{\partial}{\partial x^i}.
\]
向量分量的逆变（上标）变换为
\[
\bar v^a=\frac{\partial \bar x^a}{\partial x^i}v^i,
\]
而协变（下标）分量变换相应为
\[
\bar v_a=\frac{\partial x^i}{\partial \bar x^a}v_i.
\]
度量张量分量亦按张量变换规律变换：
\[
\bar g_{ab}=\frac{\partial x^i}{\partial \bar x^a}\frac{\partial x^j}{\partial \bar x^b}g_{ij}.
\]

\subsection{备注与应用}
在处理物理问题时，选择与问题对称性相匹配的坐标系通常能简化方程与积分。例如柱对称问题常用圆柱坐标，球对称问题常用球坐标。对于一般曲面或广义坐标系，度量张量提供了完整的几何信息，用以计算面积元、体积元以及在该坐标系下的梯度、散度、旋度和拉普拉斯算子的具体表达式。

\subsection{示意图（教学用）}
下面给出三个简单示意图以辅助直观理解（如果想要更精细的图我可以改进）：
\begin{figure}[htbp]
\centering
\begin{tikzpicture}[scale=1]
	\draw[->] (0,0) -- (2,0) node[right] {$x$};
	\draw[->] (0,0) -- (0,2) node[above] {$y$};
	\draw[->] (0,0) -- (-0.8,-0.8) node[below left] {$z$};
	\node at (1.0,1.0) {直角坐标};
\end{tikzpicture}
\hspace{1cm}
\begin{tikzpicture}[scale=0.9]
	\draw[->] (0,0) -- (1.8,0) node[right] {$r$};
	\draw (0,0) circle (1.2);
	\draw[->] (0.8,0) arc (0:40:0.8) node[midway, right] {$\phi$};
	\node at (0.5,-1.4) {柱坐标平面};
\end{tikzpicture}
\hspace{1cm}
\begin{tikzpicture}[scale=0.9]
	\draw[->] (0,0) -- (1.6,0) node[right] {$r\sin\theta\cos\phi$};
	\draw (0,0) -- (0.6,1.4) node[above right] {$\theta$};
	\draw (0,0) -- (-0.7,0.9) node[left] {$\phi$};
	\node at (0.3,-1.2) {球坐标示意};
\end{tikzpicture}
\caption{左：直角坐标；中：柱坐标 $(r,\phi)$ 平面；右：球坐标示意}
\end{figure}

% End of coordinate section translation
\subsection{面积元与体积元}
在给定坐标系中，面积元或体积元可以由度量张量的行列式给出。对于 $n$ 维情形，体积元为
\[
\mathrm{d}V=\sqrt{|g|}\,\mathrm{d}x^{1}\cdots\mathrm{d}x^{n},
\]
其中 $g=\det(g_{ij})$。在三维正交坐标系（尺度因子 $h_1,h_2,h_3$）中，体积元简化为
\[
\mathrm{d}V=h_1h_2h_3\,\mathrm{d}u^1\mathrm{d}u^2\mathrm{d}u^3.
\]
例如：
- 柱坐标：$\mathrm{d}V=r\,\mathrm{d}r\mathrm{d}\phi\mathrm{d}z$。
- 球坐标：$\mathrm{d}V=r^2\sin\theta\,\mathrm{d}r\mathrm{d}\theta\mathrm{d}\phi$。

\subsection{常用微分算子在正交坐标系下的形式}
在尺度因子为 $h_1,h_2,h_3$ 的正交坐标系 $(u^1,u^2,u^3)$ 中，常用算子的表达式为：
\begin{itemize}
\item 梯度（标量场 $\Phi$）：
\[
\nabla\Phi=\sum_{i=1}^3 \frac{1}{h_i}\frac{\partial\Phi}{\partial u^i}\hat{\mathbf{e}}_i.
\]
\item 散度（向量场 $\mathbf{A}=A_i\,\hat{\mathbf{e}}_i$）：
\[
\nabla\cdot\mathbf{A}=\frac{1}{h_1h_2h_3}\sum_{i=1}^3 \frac{\partial}{\partial u^i}\bigl(h_1h_2h_3\frac{A_i}{h_i}\bigr).
\]
\item 拉普拉斯算子（标量）：
\[
\nabla^2\Phi=\frac{1}{h_1h_2h_3}\sum_{i=1}^3 \frac{\partial}{\partial u^i}\Bigl(\frac{h_1h_2h_3}{h_i^2}\frac{\partial\Phi}{\partial u^i}\Bigr).
\]
\end{itemize}

\subsection{示例：球坐标中径向场的散度}
考虑向量场 $\mathbf{A}=A_r(r)\,\hat{\mathbf{e}}_r$（仅径向分量，与角无关），在球坐标中 $h_r=1,h_\theta=r,h_\phi=r\sin\theta$，由散度公式得：
\[
\nabla\cdot\mathbf{A}=\frac{1}{r^2}\frac{\partial}{\partial r}(r^2 A_r(r)).
\]
若 $A_r=\dfrac{C}{r^2}$（例如点源势场的径向通量），则
\[
\nabla\cdot\mathbf{A}=\frac{1}{r^2}\frac{\partial}{\partial r}(r^2\cdot \frac{C}{r^2})=0\quad(r\neq0),
\]
表明在原点外该场源项为零。

\subsection{TikZ 示意图}
% 下面给出三个简单的 TikZ 图示：直角坐标系、柱坐标 (r,phi) 平面、球坐标示意
\begin{figure}[htbp]
\centering
\begin{tikzpicture}[scale=1]
	% Cartesian axes
	\draw[->] (0,0) -- (2,0) node[right] {$x$};
	\draw[->] (0,0) -- (0,2) node[above] {$y$};
	\draw[->] (0,0) -- (-0.8,-0.8) node[below left] {$z$};
	\node at (1.0,1.0) {直角坐标};
\end{tikzpicture}
\hspace{1cm}
\begin{tikzpicture}[scale=0.9]
	% Cylindrical r-phi plane
	\draw[->] (0,0) -- (1.8,0) node[right] {$r$};
	\draw (0,0) circle (1.2);
	\draw[->] (0.8,0) arc (0:40:0.8) node[midway, right] {$\phi$};
	\node at (0.5,-1.4) {柱坐标平面};
\end{tikzpicture}
\hspace{1cm}
\begin{tikzpicture}[scale=0.9]
	% Spherical radius and angles
	\draw[->] (0,0) -- (1.6,0) node[right] {$r\sin\theta\cos\phi$};
	\draw (0,0) -- (0.6,1.4) node[above right] {$\theta$};
	\draw (0,0) -- (-0.7,0.9) node[left] {$\phi$};
	\node at (0.3,-1.2) {球坐标示意};
\end{tikzpicture}
\caption{左：直角坐标；中：柱坐标 $(r,\phi)$ 平面；右：球坐标示意}
\end{figure}

% End of additions
% End of file
